\section{Constructing stable function domains and their PCF model}
\label{sec:p10-stable-domains-sequential-models}

The continuous PCF model in Section~\ref{sec:p10-operational-adequacy-domain-models}
uses full function spaces with pointwise order. Gunter's other interpretation
uses stable maps with a different order; its existence is not established by
changing the word ``continuous'' to ``stable''. We construct the stable
function domain, its finite approximants and its evaluation/recursion maps
\cite[Definitions~4.4.3--4; Lemma~4.4.5, pp.~35--36]{Gunter1994TypesAuthorOA}.
The modulus interpretation is from Abramsky--Jung's corrected/expanded
author chapter \cite[Section~8.2, p.~147]{AbramskyJungCorrectedExpandedDomainTheory}.
The construction here concerns pointed dI-domains and unrestricted stable
maps, not a claim that every stable domain theory or bidomain presentation
has these same hypotheses.

\begin{hbconvention}{def:p10-sd-domains}{Distributivity, finite compact ideals and compatibility}
A dI-domain here is a set-sized pointed algebraic dcpo, bounded-complete,
distributive, with property I: \(\downarrow a\) is finite for each
compact \(a\in K(D)\). Bounded-completeness supplies joins of all
bounded subsets; the empty join is bottom. Write \(x\mathrel{\uparrow}y\)
when \(x,y\) have a common upper bound. Nonempty infima exist, and
distributivity means
\[
 x\sqcap(y\sqcup z)=(x\sqcap y)\sqcup(x\sqcap z)
                       \qquad(y\mathrel{\uparrow}z).
\]
The compatibility condition is on \(y,z\), whose join is being used;
it is not replaced by compatibility of \(x,y\).
All maps called stable below are Scott-continuous and preserve meets of
compatible pairs. They need not be strict: constant maps are included.
\cite[pp.~35--36]{Gunter1994TypesAuthorOA}
\cite[Proposition~4.1.2, p.~56]{AbramskyJungCorrectedExpandedDomainTheory}.
\end{hbconvention}

\begin{hbtheorem}{thm:p10-sd-foundations}{proved=annex:p10.sd.foundations}{The finite-ideal facts needed by the construction}
Every element below a compact element of a dI-domain is compact.
Finite bounded joins of compact elements are compact. Binary meet is
jointly Scott-continuous. Every principal ideal is a distributive complete
lattice. Finite products of dI-domains are dI-domains, with coordinatewise
order and operations. These statements include the one-point product.
\end{hbtheorem}

\begin{hbdefinition}{def:p10-sd-stable-order}{Stable functions, moduli and their stronger order}
For dI-domains \(D,E\), let \([D\longrightarrow_s E]\) consist of
Scott-continuous functions satisfying
\(f(x\sqcap y)=f(x)\sqcap f(y)\) whenever \(x\mathrel{\uparrow}y\).
Its source-native stable order is
\[
 f\sqsubseteq_s g\quad\Longleftrightarrow\quad
 (\forall x\sqsubseteq y)\ f(x)=f(y)\sqcap g(x).
\]
Ordinary pointwise order is denoted \(f\sqsubseteq g\).
For compact output information \(b\sqsubseteq f(x)\), the modulus
\(M(f,x,b)\) is the least \(d\sqsubseteq x\) with
\(b\sqsubseteq f(d)\), when it exists. Thus the output must contain
the required information, not be below it. This is the compact algebraic
case of the source's minimum-input notion
\cite[p.~36]{Gunter1994TypesAuthorOA}
\cite[p.~147]{AbramskyJungCorrectedExpandedDomainTheory}.
\end{hbdefinition}

\begin{hbtheorem}{thm:p10-sd-moduli}{proved=annex:p10.sd.moduli}{Actual minimum witnesses and the stable-order equivalence}
A continuous map between dI-domains is stable iff these moduli exist for
every \(x\) and compact \(b\sqsubseteq f(x)\). Each modulus is compact.
Such a map preserves every bounded nonempty infimum.
For stable maps,
\[
 f\sqsubseteq_s g\quad\Longleftrightarrow\quad
 f\sqsubseteq g\ \text{and}\
 M(f,x,b)=M(g,x,b)\quad(b\in K(E),\ b\sqsubseteq f(x)).
\]
The stable order therefore records required input information, not merely
pointwise output refinement.
\end{hbtheorem}

\begin{hbtheorem}{thm:p10-sd-ordered-functions}{proved=annex:p10.sd.ordered-functions}{The actual stable function dcpo and its bounded joins}
The stable order is a partial order with the constant-bottom map least.
Directed suprema are pointwise suprema and remain stable maps.
Every stable-order bounded family has its stable supremum pointwise.
If \(f,g\sqsubseteq_s h\), then their stable meet is pointwise meet.
Within \(\downarrow_s h\), stable order agrees with pointwise order;
this need not hold for an arbitrary pair of stable functions.
\end{hbtheorem}

\begin{hbdefinition}{def:p10-sd-finite-cuts}{Finite cut functions, not an assumed function-space basis}
For stable \(f:D\to E\), compact \(a\in K(D)\), \(b\in K(E)\), put
\[
                     f_{a,b}(x)=f(x\sqcap a)\sqcap b.
\]
Its input cut lies in the finite ideal \(\downarrow a\), and its
range in \(\downarrow b\). A finite join of such cuts is formed only
when they have a common stable upper bound, for example when all are
cuts of the same \(f\). No unbounded pointwise join is stipulated.
\end{hbdefinition}

\begin{hbtheorem}{thm:p10-sd-function-domain}{proved=annex:p10.sd.function-domain}{Finite approximants prove full dI closure}
Each cut \(f_{a,b}\) is stable, compact in the stable function dcpo,
and \(f_{a,b}\sqsubseteq_s f\). The directed family of finite joins
of cuts of \(f\), including the empty bottom join, has supremum \(f\).
Every compact stable function is such a finite join of its own cuts.
Its stable principal ideal is finite. The function dcpo is distributive
and bounded-complete, hence a dI-domain.
Countable compact bases in \(D,E\) give a countable compact basis in
\([D\longrightarrow_s E]\). This proves the source's asserted
closure rather than using it as a premise
\cite[p.~36]{Gunter1994TypesAuthorOA}.
\end{hbtheorem}

\begin{hbtheorem}{thm:p10-sd-cartesian-closed}{proved=annex:p10.sd.cartesian-closed}{Stable evaluation, currying and composition}
Pointed dI-domains and all stable maps form a cartesian closed category.
Products are the actual coordinate products and the exponential is the
stable function domain. Evaluation \((f,x)\mapsto f(x)\) is jointly
continuous and stable. Currying gives inverse bijections of stable maps
\[
 [C\times D\longrightarrow_s E]
       \cong[C\longrightarrow_s[D\longrightarrow_s E]].
\]
Composition is monotone and preserves directed suprema in the displayed
stable hom orders. None of these morphisms is assumed strict.
\end{hbtheorem}

\begin{hbtheorem}{thm:p10-sd-stable-fix}{proved=annex:p10.sd.stable-fix}{The stable fixed-point operator is proved, not postulated}
For a dI-domain \(D\), the ordinary least fixed point of a stable
endomap gives a Scott-continuous stable map
\[
 \operatorname{fix}:[D\longrightarrow_s D]\longrightarrow_s D.
\]
For compatible endomaps with common stable upper bound,
\[
 \operatorname{fix}(f\sqcap_s g)
       =\operatorname{fix}(f)\sqcap\operatorname{fix}(g).
\]
This assertion is about the stable hom order; its proof tracks both
iteration chains under their common upper-bound chain.
\end{hbtheorem}

\begin{hbdefinition}{def:p10-sd-pcf-model}{Gunter's stable interpretation of the already specified PCF}
Keep exactly the language/evaluation of Definitions~\ref{def:p10-oa-language}
and~\ref{def:p10-oa-evaluation}, including \(\operatorname{pred}(0)=0\).
Write \(S_t=D\llbracket t\rrbracket\) for the source's stable
interpretation, distinct from its continuous interpretation \(C\):
\[
 S_{\mathsf{num}}=\mathbb N_{\bot},\quad
 S_{\mathsf{bool}}=\{\bot,\mathsf{true},\mathsf{false}\},\quad
 S_{s\to t}=[S_s\longrightarrow_s S_t].
\]
Use the stable evaluation/currying just constructed; ground constructors
and conditional retain their specified values; interpret \(\mu\) by
the proved stable fixed-point operator. No PCF parallel primitive is added.
The source calls the boolean carrier \(T=S_{\mathsf{bool}}\); here
\(T\) is a domain, not the earlier closed PCF counterexample term.
\cite[Lemma~4.4.5, p.~36]{Gunter1994TypesAuthorOA}.
\end{hbdefinition}

\begin{hbtheorem}{thm:p10-sd-pcf-interpretation}{proved=annex:p10.sd.pcf-interpretation}{Stable environment dependence and sound substitution}
Every stable PCF interpretation is a stable map from its finite environment
product to its result domain. It respects bound renaming and typed
substitution, and each registered call-by-name contraction preserves it.
In particular the body of a typed recursion binder induces an actual
stable endomap and the binder's meaning is well defined.
\end{hbtheorem}

\begin{hbdefinition}{def:p10-sd-relation}{The adequacy relation on stable type carriers}
For \(a\in S_t\) and a closed \(M:t\), write
\(a\mathrel{\mathcal R^s_t}M\). At ground type bottom is related
to every term, and a nonbottom point is related exactly when the term
evaluates to its corresponding numeral or boolean. At an arrow type,
\[
 f\mathrel{\mathcal R^s_{u\to t}}M\quad\Longleftrightarrow\quad
 (\forall N:u\text{ closed})(\forall a\in S_u)\
 a\mathrel{\mathcal R^s_u}N\ \Longrightarrow\
 f(a)\mathrel{\mathcal R^s_t}(M N).
\]
The arrow element is in the stable function carrier. This definition
retains the old operational language but changes the semantic carriers;
its closure and fundamental lemma are proved below.
\end{hbdefinition}

\begin{hbtheorem}{thm:p10-sd-pcf-adequacy}{proved=annex:p10.sd.pcf-adequacy}{Both stable-model observation directions and contextual adequacy}
For closed ground programs and ground values,
\[
 M\Downarrow V\quad\Longleftrightarrow\quad
             D\llbracket M\rrbracket=D\llbracket V\rrbracket,
 \qquad
 M\Uparrow\quad\Longleftrightarrow\quad D\llbracket M\rrbracket=\bot.
\]
Stable-model denotational equality implies contextual operational
equivalence. The typed semantic-to-syntax logical relation is instantiated
on \(S_t\), not assumed equal to the earlier relation on the continuous
type carriers. Its recursion and directed-closure dependencies are checked.
\end{hbtheorem}

\begin{hbtheorem}{thm:p10-sd-counterexample}{proved=annex:p10.sd.counterexample}{Pointwise/stable incomparability and the source's second failure of full abstraction}
On the flat boolean domain let
\(p(x)=\mathsf{true}\) and
\(q(\bot)=\bot\), \(q(\mathsf{true})=q(\mathsf{false})=\mathsf{true}\).
Both are stable; \(q\sqsubseteq p\) pointwise, but they are incomparable
and incompatible in the stable order. Their stable meet is the bottom
function although their pointwise meet is \(q\). On the stable boolean function
domain define
\[
 r(h)=\begin{cases}\mathsf{true},&h=q,\\
                    \mathsf{false},&h=p,\\\bot,&\text{otherwise}.
       \end{cases}
\]
This is continuous and stable. The functions \(p,q\) are represented
by the closed PCF abbreviations
\(p=\lambda x:\mathsf{bool}.\mathsf{true}\) and
\(q=\lambda x:\mathsf{bool}.\mathsf{if}\ x\ \mathsf{then}\ \mathsf{true}\ \mathsf{else}\ \mathsf{true}\).
With \(\Omega=\Omega_{\mathsf{bool}}\), set
\[
 \begin{split}
 M={}&\lambda f:(\mathsf{bool}\to\mathsf{bool})\to\mathsf{bool}.
    \mathsf{if}\ f q\ \mathsf{then}\
    (\mathsf{if}\ f p\ \mathsf{then}\ \Omega\ \mathsf{else}\ \mathsf{true})
       \ \mathsf{else}\ \Omega,\\
 N={}&\lambda f:(\mathsf{bool}\to\mathsf{bool})\to\mathsf{bool}.
    \mathsf{if}\ f q\ \mathsf{then}\
    (\mathsf{if}\ f p\ \mathsf{then}\ \Omega\ \mathsf{else}\ \mathsf{false})
       \ \mathsf{else}\ \Omega.
 \end{split}
\]
The source-native names \(p,q,r,M,N\) denote these typed functions
and abbreviations, not ambient predicate or set variables.
Then \(M\simeq_o N\), while
\(D\llbracket M\rrbracket(r)=\mathsf{true}\ne\mathsf{false}=D\llbracket N\rrbracket(r)\).
Thus the stable model is adequate but not fully abstract. Parallel-or from
the previous counterexample is itself not stable. These are different
failures: \(r\) is not the denotation of any closed term of this PCF.
Stable maps therefore do not coincide with its sequentially definable programs
\cite[p.~38]{Gunter1994TypesAuthorOA}.
\end{hbtheorem}

\begin{proofboundary}{def:p10-sd-scope}{A stable construction is not a bidomain classification}
The proved construction is the pointed dI stable-function model, with
finite-ideal hypotheses used explicitly. General continuous stable
L-domains, concrete-domain sequential algorithms, two-order bidomain
axioms/classifications, and stable recursive type equations have separate
construction obligations. Merely displaying both pointwise and stable
orders does not establish a bidomain theorem. The distinct Berry/Girard/
coherence-space presentations referred to by the sources, their original
papers, and complete printed/author manifestations remain separate duties.
\end{proofboundary}
